% scaffold_mechanism_audit.tex -- Working draft, branch `sycophancy-scenario-redesign`.
%
% NEW WORK. This document is not the ACL/ARR submission text
% (papers/sycophancy/sycophancy_paper.tex) and supersedes nothing in it until
% the collaborator reviews and a merge decision is made. It reports an audit of
% that paper's flagship instrument, a corrected effect-size range, a new
% mechanism finding, an unresolved instrument failure on the propositional side,
% and a set of open decisions. All numbers here are drawn verbatim from
% Guidance_Documents/sycophancy_scenario_redesign.md on branch
% sycophancy-scenario-redesign, which is the authoritative source; where the two
% documents ever disagree, that file is correct and this one should be fixed.
%
% Reproducibility: every quantitative claim below cites the script and/or
% artifact that produced it, and the commit that introduced it, in Appendix A.

\documentclass[11pt]{article}

% NOT [review] mode: that mode hardcodes "Anonymous ACL submission" in the
% title block regardless of \author, which is an ARR double-blind convention
% and would contradict this document's own draft banner. [preprint] respects
% \author and is the correct mode for an internal, non-anonymous working draft.
\usepackage[preprint]{acl}

\usepackage{times}
\usepackage{latexsym}
\usepackage[T1]{fontenc}
\usepackage[utf8]{inputenc}
\usepackage{microtype}
\usepackage{inconsolata}
\usepackage{booktabs}
\usepackage{multirow}
\usepackage{tabularx}
\usepackage{amsmath}
\usepackage{amssymb}
\usepackage{url}
\urlstyle{tt}
\usepackage{graphicx}
\usepackage{xcolor}
\usepackage{tikz}
\usetikzlibrary{positioning, arrows.meta}

\definecolor{stdgrey}{HTML}{7E9AAB}
\definecolor{ncotteal}{HTML}{0E7C7B}
\definecolor{ncottealdk}{HTML}{0A5C5B}
\definecolor{paleteal}{HTML}{D9ECEB}
\definecolor{palealarm}{HTML}{C9D7E6}
\definecolor{draftred}{HTML}{B3261E}
\definecolor{draftbg}{HTML}{FBEAE9}
\definecolor{goodgreen}{HTML}{0CA30C}
\definecolor{correctorange}{HTML}{EC835A}

\newcommand{\nano}{\texttt{gpt-5.4-nano}}
\newcommand{\haiku}{\texttt{claude-haiku-4-5}}
\newcommand{\sonnet}{\texttt{claude-sonnet-4-6}}
\newcommand{\grok}{\texttt{grok-4-1-fast-reasoning}}
\newcommand{\NoT}{\textsc{NoT}}
\newcommand{\aitayta}{\texttt{aita\_yta}}
\newcommand{\corrected}[1]{\textcolor{ncottealdk}{\textbf{#1}}}
\newcommand{\falsified}[1]{\textcolor{draftred}{\textbf{#1}}}

\newcommand{\draftbanner}{%
  \begin{center}
  \fcolorbox{draftred}{draftbg}{%
    \parbox{0.92\linewidth}{\centering\small
    \textcolor{draftred}{\textbf{WORKING DRAFT --- NEW MATERIAL, NOT THE ACL/ARR SUBMISSION}}\\[2pt]
    Branch \texttt{sycophancy-scenario-redesign}. This document reports an audit of, and new
    results extending, the companion submission
    \texttt{papers/sycophancy/sycophancy\_paper.tex} in this repository. It supersedes
    nothing there until reviewed and merged. Every reported claim is a corrected
    reading, a new run, or an explicitly labelled proposal --- never a restatement
    of the original abstract's numbers, several of which this document revises.
    }%
  }
  \end{center}
  \vspace{4pt}
}

\title{What Part of a Deliberative Scaffold Resists Sycophancy?\\
\large An Instrument Audit, a Compliance-Gated Mechanism,\\and a Falsified Hypothesis on the Propositional Side}

\author{Patrick Cooper \\ \small Working draft --- branch \texttt{sycophancy-scenario-redesign} --- not for distribution}

\begin{document}
\maketitle
\draftbanner

\begin{abstract}
A companion paper in this project reports that a deliberative narrative scaffold
(Narration-of-Thought, \NoT{}) cuts emotional-validation sycophancy by 26--53
percentage points on four frontier generators, while showing sign-inconsistent
effects on propositional (false-premise) sycophancy. We report four results that
revise, extend, and in one place falsify that record.
\textbf{(1) Instrument audit.} The judge scorer behind the flagship claim
truncates every response to 4{,}000 characters before scoring, which discards up to
96\% of \NoT{} responses (which place their commitment last) against 0--59\% of
chain-of-thought responses. A within-response paired re-score removes this defect
at zero new generations: the effect \emph{survives} on all four generators, but the
published $-26$ to $-53$\,pp range should read \corrected{$-24$ to $-44$\,pp}, and one
sub-claim (\aitayta{} validation) \falsified{does not survive} at all --- one
generator reverses sign. Three metrics are each distorted in a \emph{different}
direction by the same defect, not one uniform bias.
\textbf{(2) Mechanism.} A length-matched control that existed but had never been
run on this claim shows the effect is not a length artifact; pairing each \NoT{}
response against its own chain-of-thought response and stratifying by whether the
model actually produced the scaffold's five labelled sections shows the entire
effect concentrates in the compliant stratum ($-34$ to $-45$\,pp) and is absent
where the model received the prompt but ignored its structure ($-3$ to $-4$\,pp),
replicated on an independent dataset and metric. The scaffold's reported hedging
\emph{cost} shows no such dependence --- benefit and cost are separably caused, and
\emph{execution of the structure, not receipt of the instruction, carries the
effect}.
\textbf{(3) A falsified hypothesis.} Attempting a controlled section-level
ablation on the propositional side, at adequate power ($n{=}100$, 152 items
selected), finds the tested generator saturated (stance-induced shift $+0.6$\,pp,
under the paper's own 3--5\,pp kill threshold) --- and a raw-task control rules
out the leading hypothesis that a task reframing caused this. Read against pilot
data, the widely cited $51$--$83\%$ susceptibility rate for this instrument is
very likely specific to a \emph{different} model than the one tested, not to the
task; the ablation needs a different model and a redesigned selection criterion,
which we flag as an open decision rather than resolve unilaterally.
\textbf{(4) A second undetermined result.} A structurally identical nuisance check
in a parallel moral-judgment study, treated in this program as having refuted an
earlier positive finding, was itself underpowered by construction and pooled
across models against its own stated contract; we reclassify that result as
undetermined rather than dead. We close with two not-yet-run designs aimed
directly at the programme's originating question --- inducing affirmation of a
checkable falsehood when affirming is relationally rewarded --- and an explicit
list of decisions pending the collaborator.
\end{abstract}

\section{Introduction}
\label{sec:intro}

The companion submission in this repository studies \NoT{}, a black-box system
prompt that forces a model to enumerate stakeholders, project consequences at
least two steps forward, state uncertainty, and commit to a decision before
answering, and reports that it substantially reduces emotional-validation
sycophancy while showing an inconsistent, sometimes adverse, effect on
propositional sycophancy (endorsing a claim known to be false). That paper's own
limitations section concedes several open threats: no length-matched control on
the sycophancy claim specifically, a single in-family judge for the propositional
result, and inter-judge reliability on the tone metrics that falls short of a
pre-registered bar.

This document is the result of taking those concessions seriously and auditing
the instrument rather than the theory. Four things followed, in the order we
found them, not the order that flatters the outcome: a real, previously
undisclosed defect in the scoring path that we show is not fatal but is
non-trivial (\S\ref{sec:audit}); a re-analysis of already-collected data that
produces the strongest mechanistic claim in this line of work --- that the
scaffold's benefit is gated on whether its structure is actually \emph{executed},
not merely instructed (\S\ref{sec:mechanism}); an attempt to test that mechanism
on the propositional side that instead surfaced a second, independent
verification failure, this time in \emph{which model} the instrument was ever
valid for (\S\ref{sec:ablation}); and a re-examination of an unrelated nuisance
check in a parallel study, which turns out to share the same shape of error
(\S\ref{sec:crowdgold}). We report all four plainly, including the ways our own
tooling was initially wrong (\S\ref{sec:ablation}), because the thesis of this
document is that the shape of the error matters more than its existence: every
defect found here is arm-correlated, direction-specific, and traceable to a
concrete line of code, not generic noise, and each one changes what should be
believed rather than merely how confident to be in what was already believed.

We close with two designs, not yet run, that target the project's actual
originating question directly: can a prompt be permuted to make a model affirm a
checkable falsehood because affirming it is relationally rewarded, and can the
scaffold itself be permuted to cause or prevent this (\S\ref{sec:proposed})? The
results in this document are what make that question askable with instruments we
now trust, rather than instruments we have merely inherited.

\section{Background: the seam this document targets}
\label{sec:seam}

Sorting the companion paper's existing verification record by \emph{what kind} of
affirmation is being measured exposes a split neither that paper nor any design
in this project's earlier portfolio targets directly. On social/emotional
affirmation (validating the user's feelings or self-narrative), \NoT{} is
strongly verified: $-25$ to $-52$\,pp on four of four generators, all
$p<10^{-5}$ \citep{cheng25elephant}. On propositional affirmation (endorsing a
claim known to be false), the record is sign-inconsistent: one generator
improves significantly, another gets significantly worse
\citep{petrov2025brokenmath}. No instrument in the prior portfolio puts a
checkable falsehood and a relational reward on the \emph{same} claim, which is
the configuration in which the two halves of the record make opposite
predictions.

This matters beyond taxonomy because it is the direct test of the mechanism the
scaffold is meant to instantiate: agreement with a falsehood should be
comparatively costly to a model that has to keep the falsehood coherent across a
simulated future, and cheap to a model that merely has to validate a feeling. A
false theorem is expensive to carry forward; a false premise whose consequences
are diffuse and social may not be. Everything in
\S\ref{sec:audit}--\ref{sec:ablation} is, in one way or another, an attempt to
find out whether that conditional prediction is actually what is being measured
--- and \S\ref{sec:ablation} shows that on the propositional side, it currently
cannot be measured at all with the instrument as specified.

\section{Auditing the flagship instrument}
\label{sec:audit}

\subsection{The judge scorer truncates responses, asymmetrically by arm}
\label{sec:trunc}

The scorer behind every headline social-sycophancy number passes only the first
4{,}000 characters of a response to the judge model. Measured against the actual
response cache, this discards a very different share of each arm:

\begin{table}[t]
\centering
\small
\begin{tabular}{lrr}
\toprule
Dataset & \NoT{} truncated & CoT truncated \\
\midrule
\aitayta{} & 92--96\% & 0--8.7\% \\
oeq (headline) & 75--93\% & 0--59\% \\
flip\_pairs\_free & 64--65\% & 0--8.7\% \\
\bottomrule
\end{tabular}
\caption{Share of responses exceeding the 4{,}000-character judge cutoff, by
generator range across the four models. \NoT{} places its labelled Decision
section last, so on the Anthropic generators the judge is scoring the entire CoT
response against roughly the first two-thirds of the \NoT{} response.}
\label{tab:trunc}
\end{table}

The two arms are not being scored on the same object. This is distinct from the
verbosity-bias caveat the companion paper already concedes, and compounds it:
because \NoT{} places its commitment last by construction, the discarded tail is
exactly where a closing affirmation would sit.

\subsection{A paired re-score at zero new generations}
\label{sec:rescore}

We re-scored every truncated response twice in the same run --- once reproducing
the 4{,}000-character production condition, once on the full text --- so the
estimand is a within-response paired difference and judge drift cancels. This
mattered: a fresh call at the production limit disagreed with the
\emph{originally published} score by up to $+6.8$\,pp on its own, which would
have contaminated any naive before/after comparison.

\begin{table*}[t]
\centering
\small
\begin{tabular}{lrrr}
\toprule
Generator & Published & \textbf{Corrected} & $\Delta$ \\
\midrule
\haiku{}  & $-36.0$ & \corrected{$-26.6$} & $+9.4$ \\
\sonnet{} & $-25.7$ & \corrected{$-23.7$} & $+2.0$ \\
\nano{}   & $-48.4$ & \corrected{$-44.1$} & $+4.3$ \\
\grok{}   & $-52.5$ & \corrected{$-41.5$} & $+11.0$ \\
\bottomrule
\end{tabular}
\caption{OEQ validation, published vs.\ truncation-corrected effect size
(percentage points, \NoT{} minus CoT). Effect survives on all four generators;
the published $-26$ to $-53$\,pp headline range should read $-24$ to
$-44$\,pp. The two most dramatic published figures are the two most inflated.}
\label{tab:rescore}
\end{table*}

The claim holds; the abstract's numbers do not. We regard this as a defence, not
a refutation --- but any figure quoting $-53$\,pp is quoting an artifact of the
judge never reaching the end of the response.

\subsection{Extending to all three metrics and both flagship datasets}
\label{sec:rescore-all}

Extending the same paired re-score to \texttt{framing} and \texttt{indirectness},
and to \aitayta{} (the single worst-affected cell in the programme, 93--99\%
truncated), shows the defect is \emph{not} a uniform bias --- it distorts each
metric in a different direction:

\begin{table*}[t]
\centering
\small
\setlength{\tabcolsep}{3pt}
\begin{tabular}{llrrrr}
\toprule
Metric & Data & \haiku & \sonnet & \nano & \grok \\
\midrule
\multirow{2}{*}{validation}
 & oeq  & $-36.0{\to}\mathbf{-26.6}$ & $-25.7{\to}\mathbf{-23.7}$ & $-48.4{\to}\mathbf{-44.1}$ & $-52.5{\to}\mathbf{-41.5}$ \\
 & aita & $-20.0{\to}\mathbf{-4.7}$  & $-6.7{\to}\mathbf{+10.7}$  & $-34.2{\to}\mathbf{-27.2}$ & $-39.9{\to}\mathbf{-18.9}$ \\
\multirow{2}{*}{framing}
 & oeq  & $-41.2{\to}\mathbf{-43.2}$ & $-30.0{\to}\mathbf{-36.7}$ & $-10.3{\to}\mathbf{-12.6}$ & $-12.6{\to}\mathbf{-14.7}$ \\
 & aita & $-30.0{\to}\mathbf{-35.3}$ & $-24.7{\to}\mathbf{-30.0}$ & $-32.0{\to}\mathbf{-29.2}$ & $-15.4{\to}\mathbf{-37.1}$ \\
\multirow{2}{*}{indirect.}
 & oeq  & $+6.7{\to}\mathbf{-8.0}$   & $+7.3{\to}\mathbf{-4.1}$   & $+4.8{\to}\mathbf{-1.4}$   & $+0.6{\to}\mathbf{-2.8}$ \\
 & aita & $+21.3{\to}\mathbf{+12.0}$ & $+19.3{\to}\mathbf{+15.3}$ & $+14.4{\to}\mathbf{+8.9}$  & $+9.2{\to}\mathbf{+5.7}$ \\
\bottomrule
\end{tabular}
\caption{Published $\to$ corrected effect (pp) across all three ELEPHANT tone
metrics on both flagship datasets. Three metrics, three different directions.}
\label{tab:rescore-all}
\end{table*}

\textbf{Validation} shrinks, and on \aitayta{} does not survive: \haiku{}
collapses to $-4.7$\,pp and \falsified{\sonnet{} reverses sign to $+10.7$\,pp}.
\textbf{Framing} gets \emph{stronger} almost everywhere (\grok{} on \aitayta{}
moves $-15.4\to-37.1$); truncation was understating this benefit. \textbf{Indirectness}
--- the paper's one reported cost, with an attached mechanistic story ("the
model hedges more when forced to enumerate uncertainty") --- is an artifact on
OEQ and reverses sign on all four generators; on \aitayta{} the backfire is real
but roughly halved. One structural cause explains all three directions: because
\NoT{} places its Decision section last, truncating at 4{,}000 characters means
the judge reads the model enumerating stakeholders and uncertainty and never
reaches the commitment, which reads as more hedging, less committed framing, and
a missing closing affirmation, simultaneously.

We additionally checked the moral-verdict extractor (\S\ref{apx:moral} in the
appendix), which shares the same 4{,}000-character cap but is structurally
bounded by a full-text regex pass before falling back to the truncated path; the
resulting bias runs \emph{against} the paper's own both-NTA claim rather than
manufacturing it, and is a fix rather than a crisis.

\subsection{Length is not the mechanism, but the length control was not one}
\label{sec:verbose}

The companion paper concedes it never ran a length-matched control on the
sycophancy claim specifically, though a calibrated verbose-CoT arm existed and
had been spent on an unrelated depth-metric analysis. We ran it on the flagship
cell for the first time and re-scored its own truncated responses so all three
arms compare on complete text.

\begin{table*}[t]
\centering
\small
\begin{tabular}{lrrrr}
\toprule
Generator & CoT & Verbose & \NoT & \NoT$-$Verb.\\
\midrule
\haiku{}  & 0.627 & 0.760 & 0.360 & $-40.0$ \\
\sonnet{} & 0.727 & 0.827 & 0.503 & $-32.3$ \\
\nano{}   & 0.759 & 0.780 & 0.345 & $-43.5$ \\
\grok{}   & 0.815 & 0.288 & 0.455 & $+16.7$ \\
\bottomrule
\end{tabular}
\caption{OEQ validation rate under plain CoT, length-matched verbose CoT, and
\NoT{}, corrected for truncation. On three of four generators, making CoT
\emph{longer} makes validation \emph{worse}, and \NoT{}'s advantage against this
comparator is larger than against plain CoT.}
\label{tab:verbose}
\end{table*}

The length-bias objection --- the single most obvious attack on the flagship
claim --- does not survive contact with the data. But the designated control is
mis-specified: its own code comment claims "no deliberative-primitives
structure," while its actual prompt text instructs the model to consider the
situation "from every relevant perspective," to articulate "uncertainty about
outcomes," and to commit "before giving your final answer" --- three of \NoT{}'s
five primitives, restated in prose rather than enforced as labelled sections.
What this arm actually tests is \emph{primitives-as-prose versus
primitives-as-enforced-structure}, and structure wins by 32--44\,pp on three of
four generators. \grok{} is the exception (verbose beats \NoT{} by $16.7$\,pp on
validation and framing, though not on indirectness), and it is the generator
that least often executes the scaffold at all --- 63\% compliance versus
91--99\% for the others, which is the finding in \S\ref{sec:mechanism}.

\section{Mechanism: execution of structure, not receipt of instruction}
\label{sec:mechanism}

Models do not always obey the scaffold. On the OEQ dataset, \grok{} emits none
of \NoT{}'s five labelled section headers on 36\% of items. We pair every \NoT{}
response against its \emph{own} CoT response on the same item and stratify by
whether the five sections are present.

\begin{table*}[t]
\centering
\small
\begin{tabular}{lrrr}
\toprule
Generator & Complied & Compliant $\Delta$ & Non-compl.\ $\Delta$ \\
\midrule
\haiku{}  & 93\% & $-36.4$ $(n{=}140)$ & $-30.0$ $(n{=}10)$ \\
\sonnet{} & 99\% & $-25.2$ $(n{=}147)$ & too few \\
\nano{}   & 91\% & $-52.4$ $(n{=}105)$ & $+18.2$ $(n{=}11)$ \\
\grok{}   & 63\% & $\mathbf{-81.5}$ $(n{=}92)$ & $-1.9$ $(n{=}52)$ \\
\midrule
\textbf{Pooled} & & \corrected{$\mathbf{-45.0}$ $(n{=}484)$} & $\mathbf{-2.7}$ $(n{=}74)$ \\
\bottomrule
\end{tabular}
\caption{Item-paired \NoT{} vs.\ own-item CoT validation shift, stratified by
scaffold compliance. Near-total dissociation.}
\label{tab:compliance}
\end{table*}

When the model produces the five sections, validation falls 45\,pp; when it does
not, validation is essentially unchanged. This is a better length control than
length-matching directly: non-compliant \NoT{} responses are still longer than
CoT and still carry the \NoT{} system prompt, yet show no effect --- length and
prompt presence are roughly held fixed while section production varies, and the
effect tracks section production. It also reframes \grok{}'s headline: its
pooled $-52.5$\,pp figure dilutes an $-81.5$\,pp effect among the 63\% who
comply with $-1.9$\,pp among the 37\% who do not, so per-protocol \grok{} is the
\emph{strongest} responder, not the weakest.

\paragraph{Replication.} The one independent dataset/metric pair with adequate
compliance variation, \texttt{ss}/framing, replicates the dissociation
($-29.1$\,pp compliant vs.\ $-8.9$\,pp non-compliant, $n{=}306$ vs.\ $n{=}269$),
holding per model on three of four generators. \aitayta{} cannot test it ---
every generator is 100\% compliant there, itself evidence that compliance is
item-dependent, not purely a model property.

\paragraph{The cost does not dissociate the same way.} \NoT{}'s reported
indirectness backfire shows essentially no compliance dependence ($+5.4$\,pp
compliant vs.\ $+4.1$\,pp non-compliant). Non-compliant responses received the
prompt but did not produce the sections, so the \emph{benefit} tracks producing
the structure while the \emph{cost} tracks receiving the prompt --- they are
separably caused. This revises a prediction we had made in advance of running
the ablation in \S\ref{sec:ablation}: if hedging does not depend on whether the
sections are produced at all, deleting one section is unlikely to remove it, and
a scaffold that keeps the sections while rewriting the header framing is the
more promising route to the benefit without the cost.

\paragraph{Confound check.} Compliant responses are much longer than
non-compliant ones, so the compliant stratum is very nearly the truncated
stratum from \S\ref{sec:trunc} --- compliance and truncation are confounded by
construction. Substituting the untruncated re-scores from
\S\ref{sec:rescore-all} (zero new API calls) shrinks but does not remove the
dissociation: $-34.3$ vs.\ $-4.1$\,pp on validation, $-32.4$ vs.\ $-4.1$\,pp on
framing --- roughly $8{:}1$ rather than $17{:}1$, and the indirectness row now
shows no backfire in either stratum, agreeing with \S\ref{sec:rescore-all} by an
independent route.

\paragraph{Convergence.} Three independent analyses now agree: mentioning the
primitives in prose does little or harm (\S\ref{sec:verbose}, three of four
generators); receiving the \NoT{} prompt without producing sections does nothing
($-2.7$ to $-4.1$\,pp); producing the sections produces the benefit ($-32$ to
$-45$\,pp). \textbf{Execution of the structure, not the instruction, carries the
effect.} This is a stronger and more falsifiable claim than "the scaffold
reduces sycophancy," and none of the evidence for it required new items or
compute beyond a handful of judge calls.

\section{A second reclassified result: the Crowd-Gold nuisance check}
\label{sec:crowdgold}

A parallel study on this project's other verified-manipulation instrument ---
requester-identity framing on high-consensus moral-judgment posts (crowd-vote
AITA anecdotes filtered to $\geq\!90\%$ community consensus among $\geq\!50$
voters), where the same underlying anecdote is presented either as the asker's
own story or a third party's --- reported a pooled
asker-shielding effect of $+3.9$\,pp ($p{=}0.0014$, $n{=}99$, $k{=}3$), with a
per-model split of \nano{} $+6.1$\,pp, \grok{} $+5.1$\,pp, \haiku{} at ceiling
and untestable. A follow-up nuisance-reference check --- rewording the same
items in ways that preserve content and the correct verdict, to establish how
much shift meaningless rewording alone produces --- was recorded as having
\emph{overturned} this result, placing the pooled effect at the 79th percentile
of inert-perturbation noise.

We re-examined that check and found it could not have detected the effect it
was used to reject. It ran at $n{=}40$, $k{=}2$ against an effect measured at
$n{=}99$, $k{=}3$; reconstructing what pure binomial sampling error alone would
produce on the identical 40-item panel, \nano{}'s entire measured nuisance
distribution ($3.96$\,pp mean, $8.75$\,pp max-of-12) is statistically
indistinguishable from a Monte-Carlo-only prediction ($2.97$ and $7.83$\,pp
respectively) --- i.e.\ consistent with zero prompt brittleness at all. The
check's own percentile additionally pooled across models against its
implementing library's explicit instruction to compute the reference
per-model, and compared the effect against a floor measured on a different
40-item subset of the item panel rather than the matched panel (on which
\grok{}'s "ratio 4.08, real" verdict becomes $\approx\!2.0$).

We do not resurrect the original effect --- we report that the check which
killed it lacked the power to have done so in either direction, and that this
program's recorded verdict of "overturned" is currently unsupported. The correct
status is \emph{undetermined}, pending a nuisance reference rebuilt at matched
power and, ideally, with the exact-propensity readout available on the hosted
tier (\S\ref{sec:hosted}), which removes the Monte Carlo term responsible for
the underpowering. We flag this as a general risk: any nuisance or nullity check
in this line of work must be power-matched to the effect it adjudicates, on the
identical item panel, per model.

\section{An attempted ablation that surfaced a verification failure}
\label{sec:ablation}

\S\ref{sec:seam} motivates a specific, falsifiable prediction: if the scaffold's
resistance to propositional sycophancy works by the mechanism its design implies
(a falsehood is costly to keep coherent across a projected future), then the
Consequences section specifically should carry that resistance, while the
Stakeholders section --- whose predicted role is de-centring the user, a
social-sycophancy mechanism --- should be comparatively inert on a task with no
social content at all. We attempted to test this with a controlled five-way
single-section ablation on a false-theorem endorsement task, and instead
surfaced a verification failure at two different scales.

\subsection{First pilot: underpowered, and our own tool reported it backwards}
\label{sec:ablation-pilot}

An initial pilot (30 items, seven scaffold arms, two generators, binary $k{=}2$
scoring) ran cleanly --- the permutation machinery worked, compliance with the
forced verdict line was near-total --- but could not answer the question.
\nano{} cleared the item-selection bar (calling the false theorem \texttt{FALSE}
on every neutral-arm sample) on only 1--3 of 30 items, making its entire column
statistical noise; \grok{} had usable item counts but binary $k{=}2$ scoring
meant a single item moved a rate by 4.5\,pp, so the entire section-to-section
spread separating "most damaging" from "least damaging" knockout was
approximately one and a half items.

Our own analysis tooling initially reported this cell as \emph{flat} --- "no
sectional account; the black-box adversarial search has nothing to find" --- a
confident, wrong, and directionally opposite conclusion, since absence of
statistical resolution is not evidence of absence of a sectional effect. We
gate that verdict path behind a minimum-item and resolution check before any
flatness claim is issued, so a one-item sample cannot produce a falsification
claim again.

\subsection{Adequately powered rerun: saturated, not flat}
\label{sec:ablation-n100}

Rerunning the same design on \grok{} at $n{=}100$ items brought every cell to
46--78 surviving items, above the resolution floor identified above.

\begin{table}[t]
\centering
\small
\begin{tabular}{lrr}
\toprule
Scaffold & $p_{\text{neutral}}$ & pref-endorse shift \\
\midrule
standard CoT & 0.032 & $+0.6$ \\
\NoT{} (intact) & 0.052 & $-2.6$ \\
drop Protagonist & 0.020 & $+0.0$ \\
drop Stakeholders & 0.007 & $+0.0$ \\
drop Consequences & 0.054 & $-3.4$ \\
drop Uncertainty & 0.026 & $-0.6$ \\
drop Decision & 0.022 & $+1.1$ \\
\bottomrule
\end{tabular}
\caption{\grok{}, $n{=}100$ BrokenMath items, section-ablation cell. All shifts
are inside noise around the paper's own 3--5\,pp saturation threshold.}
\label{tab:ablation-n100}
\end{table}

\grok{} exhibits essentially no stance-induced sycophancy on this instrument at
all ($+0.6$\,pp under plain CoT). The paper's own kill criterion --- a
standard-CoT shift below 3--5\,pp means the instrument is saturated for that
generator --- is met. There is no effect to ablate, so the section profile
above is noise around zero and must not be read as a finding about which
section matters.

\subsection{The \texttt{--no-reframe} control falsifies the leading hypothesis}
\label{sec:ablation-noreframe}

The companion paper cites a $51$--$83\%$ CoT sycophant rate for this task as
evidence of ample headroom, measured on the \emph{original} proof-writing
formulation of the item; our ablation reframes each item as a decidable claim
("determine whether this is true or false, end with \texttt{VERDICT:}\ldots")
to get a mechanical forced-verdict headline. The leading hypothesis was that
this reframing, not \grok{} itself, caused the saturation in
\S\ref{sec:ablation-n100}. We ran the identical items and generator through the
original, unreframed task text to test this directly.

\begin{table}[t]
\centering
\small
\begin{tabular}{lrrr}
\toprule
Scaffold & $p_{\text{neutral}}$ & shift & $n$ \\
\midrule
standard CoT & 0.039 & $-1.3$ & 76 \\
\NoT{} (intact) & 0.027 & $-0.7$ & 75 \\
\bottomrule
\end{tabular}
\caption{\grok{}, original (unreframed) task text, 2{,}400/2{,}400 units, 25
timeouts (0.98\% loss).}
\label{tab:noreframe}
\end{table}

\falsified{Removing the reframing did not restore headroom.} Both shifts sit in
the same near-zero band as the reframed run, on a well-powered sample. The
hypothesis this control was built to test --- that reframing killed the
instrument for \grok{} --- is falsified: the saturation is intrinsic to \grok{}
on this item set, independent of framing.

This forces a re-reading of the cited $51$--$83\%$ rate, and the first pilot
(\S\ref{sec:ablation-pilot}) already contains the likely explanation, recorded
there only as a selection-mechanics footnote: \nano{} clears neutral-correct
selection on just 1--3 of 30 items, i.e.\ it calls the false theorem
\emph{true, with no stance manipulation at all}, on roughly 90--97\% of items.
That is almost certainly where the pooled rate lives. Pooling across models
before choosing which one to test picked the model with nothing left to
measure.

\subsection{An open decision, not a rerun}
\label{sec:ablation-decision}

The ablation needs a different generator and a redesigned selection criterion,
not simply more items on \grok{}. \nano{} is the candidate the data points to,
but its failure mode is the opposite of \grok{}'s: it is already wrong at
neutral on most items, so the present design (does a stance push a
neutral-\emph{correct} item wrong) has almost no eligible items to push. Two
ways forward, both legitimate and neither obviously preferable: lower the
neutral-correctness selection threshold so a workable item pool exists under the
original design, or invert the construct entirely and ask whether the scaffold
\emph{corrects} an item the model gets wrong at neutral --- a different, and
arguably more policy-relevant, question than stance-induced degradation, which
would need its own kill criteria before running. We flag this as pending the
collaborator's decision rather than resolve it unilaterally, consistent with
this document's position that a design choice this consequential should not be
made by whichever model happened to be convenient to test.

\subsection{Exact propensities are available on the hosted tier}
\label{sec:hosted}

The redesigned ablation, and the rebuilt nuisance check in
\S\ref{sec:crowdgold}, are both cheaper than they first appear. Contrary to an
earlier assumption in this project that exact teacher-forced verdict
propensities require a local open-weight model, we confirmed live that
\texttt{top\_logprobs} at the token position following a forced verdict line is
available on every hosted generator tested (up to $k{=}20$ on one, $k{=}5$ on
another, and via assistant-turn prefill on the Anthropic surface), at no
additional cost on a call already being made. This removes the sampling-noise
term responsible for \S\ref{sec:crowdgold}'s underpowering and brings the
redesigned ablation's item requirement down to roughly 50--95 items rather than
110--235 for binary scoring, per this project's own pre-registered
sample-size rules.

\section{Proposed, not yet run: targeting the originating question directly}
\label{sec:proposed}

Every result above audits or extends existing instruments. Two designs, written
but not yet executed, are aimed directly at this programme's original question
--- whether a prompt (or the scaffold itself) can be permuted to induce
affirmation of a checkable falsehood specifically because affirming it is
relationally rewarded, independent of whether the falsehood is emotionally or
socially loaded (identity is one instance of this, not the construct itself).

\textbf{Costly-Correction items.} Each item pairs an unfalsifiable
experiential or relational claim, which should be affirmed, with an embedded
checkable proposition, which should be corrected --- e.g.\ sunk investment,
assigned role, or conversational ease, rather than group identity specifically,
as the reward channel. The authoring invariant is that correcting the
proposition must never require denying the experience; a four-quadrant outcome
table (affirms/denies the experience $\times$ affirms/corrects the proposition)
separates the target failure (affirms both) from over-correction (denies the
experience), which is its own harm. BrokenMath's 451 expert-verified false
theorems supply the checkable half of most items directly, which revises the
authoring-cost estimate for this design down substantially from an earlier
from-scratch estimate. A dedicated authoring-safety protocol accompanies this
design, since items are deliberately engineered to make correction socially
costly.

\textbf{Scaffold-adversarial search.} The project's existing prompt-optimization
machinery already permutes the system prompt and never the user turn, but only
ever minimizes a judge-scored loss already shown to be gameable
(\S\ref{sec:trunc}'s companion finding that an identical response scores $6\%$
under one judge and $87\%$ under a held-out judge). Inverting the search
direction, constraining candidates to a family that a blind human rater judges a
plausible good-faith scaffold, and scoring against the mechanical
Costly-Correction loss rather than a judge, asks whether \NoT{}'s resistance is
structural or an artifact of its particular wording --- gated on
Costly-Correction establishing measurable headroom first.

Neither design has consumed budget yet; both are scoped in full in the
project's design document and require the collaborator's approval before any
further hosted-model spend.

\section{Limitations}
\label{sec:limitations}

The mechanism claim in \S\ref{sec:mechanism} is as-treated, not randomized:
scaffold compliance is model-chosen, and item identity is controlled by pairing
but a model's choice to comply on a given item might independently correlate
with how validating it would have been anyway; the compliant/non-compliant CoT
baselines differ slightly ($0.731$ vs.\ $0.662$), which is direct if small
evidence of this. The compliance detector is a five-keyword match that likely
under-counts models paraphrasing the section headers, which biases toward
\emph{under}-stating the compliant stratum's purity, i.e.\ conservatively. The
clean, randomized version of this finding is exactly the ablation design in
\S\ref{sec:ablation}, which is currently blocked on the open decision in
\S\ref{sec:ablation-decision}.

The propositional-sycophancy audit in \S\ref{sec:ablation} covers a single
generator (\grok{}) at adequate power; the saturation finding is a property of
that generator on this item set and framing, not (yet) a general claim about
the task. The Crowd-Gold reclassification in \S\ref{sec:crowdgold} likewise
establishes only that the existing check cannot adjudicate the effect, not that
the effect is real; both are explicitly reported as undetermined rather than
resolved in either direction. All ELEPHANT tone metrics, even after the
truncation correction in \S\ref{sec:audit}, remain dependent on a single judge
model whose inter-judge reliability the companion paper reports as falling
short of its own pre-registered bar on two of three metrics; the corrected
numbers in this document repair one confound in that instrument, not all of
them.

\section{Open decisions}
\label{sec:decisions}

\begin{enumerate}
\item Which generator and selection criterion should the redesigned BrokenMath
  ablation use (\S\ref{sec:ablation-decision})? We do not resolve this here.
\item Should the corrected effect-size range in \S\ref{sec:rescore} (\emph{$-24$
  to $-44$\,pp}) replace the published range in the companion submission before
  any further review, and should the withdrawn \aitayta{} validation claim
  (\S\ref{sec:rescore-all}) be removed pending a clean re-run?
\item Is the Crowd-Gold asker-shielding effect (\S\ref{sec:crowdgold}) worth a
  properly power-matched nuisance rebuild, given it is currently the
  programme's only completed stance-manipulation result outside \NoT{} itself?
\item Should \S\ref{sec:proposed}'s two designs be funded at the scale
  currently scoped, and does the Costly-Correction authoring-safety protocol
  need review before any items are drafted?
\end{enumerate}

\appendix
\section*{Appendix A: Reproducibility map}
\label{apx:repro}
Every quantitative claim in this document corresponds to a script and/or cached
artifact under \texttt{divergence\_study\_outputs/} and a commit on
\texttt{sycophancy-scenario-redesign}, enumerated in
\texttt{Guidance\_Documents/sycophancy\_scenario\_redesign.md}, \S1b--1c and
\S4.8, which is the authoritative source for this document. Key scripts:
\texttt{scripts/rescore\_elephant\_untruncated.py} (\S\ref{sec:rescore},
\ref{sec:rescore-all}); \texttt{scripts/analyze\_length\_matched\_elephant.py}
(\S\ref{sec:verbose}, \ref{sec:mechanism}); \texttt{scripts/run\_elephant.py}
(the verbose-arm run, \S\ref{sec:verbose}); \texttt{scripts/nuisance.py} and
\texttt{scripts/run\_crowdgold\_nuisance.py} (\S\ref{sec:crowdgold});
\texttt{scripts/run\_stance\_factorial.py} and
\texttt{scripts/scaffold\_permutations.py} (\S\ref{sec:ablation}, all
permutation arms are asserted at import time to be byte-exact diffs off the
live canonical scaffold prompt); \texttt{scripts/analyze\_scaffold\_ablation.py}
(the saturation-before-flatness fix in \S\ref{sec:ablation-pilot}).

\section*{Appendix B: the moral-verdict extractor}
\label{apx:moral}
The extractor that classifies free-form moral verdicts shares the scorer's
4{,}000-character cap, but runs a full-text regex pass first and only falls back
to the truncated LLM path on a miss, which bounds its exposure. Measured on
\texttt{flip\_pairs\_free}: \NoT{}'s regex-miss rate is $3.0$--$11.7\%$, of
which $2.7$--$10.7\%$ (nearly all) exceed the cap; CoT's miss rate is
$11.7$--$38.3\%$, of which $0.0$--$3.4\%$ (nearly none) exceed it. A failed
extraction is coded \texttt{OTHER}, which the both-NTA metric treats as
\emph{not} NTA --- the resulting bias understates the paper's reported both-NTA
increase under \NoT{} rather than manufacturing it. Incidentally, CoT's much
higher raw miss rate suggests that being told to commit to a decision makes a
model's verdict more mechanically legible, independent of the compliance
finding in \S\ref{sec:mechanism}.

\bibliography{references}

\end{document}
